\begin{textsample}{POS Dim 1 – vem\_ed\_08 – Score 25.00 – vem\_ed\_08/t030.txt}  \label{ex:f1_pos_065}
QUEM \textbf{É} QUEM Vice-diretora- superintendente do Centro Paula Souza (CPS) desde 2019, Emilena Lorenzon Bianco \textbf{destaca}, em entrevista concedida à newsletter VEm, o valor de vivenciar projetos \textbf{colaborativos} com professores de outros países e o \textbf{que} \textbf{representa} o avanço da internacionalização do CPS.
Emilena \textbf{é} doutora em Ciência da Informação (Unesp, 2011), mestra em Engenharia de Produção (UFSCAR, 2004), especialista em Uso Estratégico das Tecnologias da Informação (Unesp, 2001) e graduada em Biblioteconomia (Unesp, 1998).
Suas pesquisas em pós-graduação abordam o cluster calçadista da região de Jaú.
Professora da Fatec Jahu entre 2004 e 2019, atuou na implantação da Fatec Matão (2018) e \textbf{foi} coordenadora de projetos na Inova Paula Souza.
Os Projetos Colaborativos Internacionais (PCIs) \textbf{são} \textbf{desenvolvidos} a distância entre professores e alunos de Fatecs e de instituições de ensino \textbf{superior} (IES) internacionais.
Qual \textbf{é} a \textbf{importância} dos PCIs na formação de \textbf{profissionais} de excelência em âmbito global? \textbf{É} extremamente importante, tanto para o aspecto \textbf{profissional}, \textbf{quanto} para a \textbf{vivência} pessoal, participar de trabalhos em parceria com outros países.
Graças à tecnologia, \textbf{que} ajudou a superar as barreiras de mobilidade física, atualmente, em \textbf{todos} os campos do \textbf{conhecimento} as práticas \textbf{são} globalizadas e os \textbf{profissionais} atuam de maneira \textbf{colaborativa}, em rede. \textbf{Programas} como os PCIs oferecem essa oportunidade para nossos professores e alunos.
Por meio da \textbf{interação} virtual, \textbf{é} \textbf{possível} oferecer escalabilidade para experimentação de outras culturas de \textbf{forma} menos custosa, o \textbf{que} \textbf{permite} incluir maior \textbf{número} de alunos.
Essas experiências agregam muito valor ao \textbf{conhecimento}.
Ao ter contato direto com \textbf{pessoas} de outras culturas, os jovens passam a entender a diversidade \textbf{profissional} nos diferentes países e como \textbf{cada} nação tem um jeito de lidar com as situações nos \textbf{ambientes} de trabalho.
Essa riqueza de visões traz inteligência cultural: a gente \textbf{aprende} como agir, como se adaptar e também desenvolve resiliência em situações desafiadoras.
Em trabalho \textbf{alinhado} com a Coordenação de Línguas da CESU, os PCIs contribuem para \textbf{promover} a Internacionalização em Casa.
Qual o \textbf{papel} dos PCIs na internacionalização dos currículos das Fatecs?
Ampliar a representatividade internacional \textbf{é} um objetivo estratégico do Centro Paula Souza.
Para \textbf{ser} \textbf{reconhecida} internacionalmente pela excelência da educação \textbf{profissional}, a instituição precisa formar \textbf{profissionais}

% matched lemmas: alinhar, ambiente, aprender, cada, colaborativo, conhecimento, desenvolvir, destacar, forma, importância, interação, número, papel, permitir, pessoa, possível, profissional, programa, promover, quanto, que, reconhecer, representar, ser, superior, todo, vivência
\end{textsample}
